\documentclass[11pt,a4paper]{article}

\usepackage[utf8]{inputenc}
\usepackage[T1]{fontenc}
\usepackage[french]{babel}
\usepackage[a4paper,margin=2cm]{geometry}
\usepackage{amsmath,amssymb}
\usepackage{lmodern}
\usepackage{enumitem}
\usepackage{array}
\usepackage{tabularx}
\usepackage{xcolor}
\usepackage{fancyhdr}
\usepackage{lastpage}
\usepackage{tikz}
\usepackage[hidelinks]{hyperref}
\usepackage{siunitx}

\setlength{\parindent}{0pt}
\setlength{\parskip}{0.4em}

\pagestyle{fancy}
\fancyhf{}
\lhead{DS de physique-chimie -- Niveau 3\ieme}
\rhead{Nom : \hspace{4cm}}
\cfoot{\thepage/\pageref{LastPage}}

\begin{document}

\begin{center}
    {\Large \textbf{DS de physique-chimie}}\\[0.3em]
    {\large \textbf{Niveau 3\ieme}}\\[0.5em]
\end{center}

\begin{center}
\renewcommand{\arraystretch}{1.4}
\begin{tabular}{|>{\centering\arraybackslash}m{3cm}|>{\centering\arraybackslash}m{3cm}|>{\centering\arraybackslash}m{6cm}|}
\hline
\textbf{Durée} & \textbf{Barème} & \textbf{Matériel autorisé} \\
\hline
1 h 30 à 2 h & 20 points & Stylo, règle, calculatrice \\
\hline
\end{tabular}
\end{center}

\textbf{Consignes générales}
\begin{itemize}[leftmargin=1.2cm]
    \item Toute réponse doit être justifiée.
    \item Les calculs doivent être rédigés avec les formules utilisées.
    \item Les unités doivent apparaître sur les résultats.
    \item Lorsqu’un document ou un rappel est donné, il peut être utilisé librement.
    \item Pensez à noter le temps que vous avez mis pour faire le sujet.
\end{itemize}

\vspace{0.3cm}

\begin{center}
\renewcommand{\arraystretch}{1.4}
\begin{tabular}{|>{\centering\arraybackslash}m{2.2cm}|>{\centering\arraybackslash}m{2.2cm}|>{\centering\arraybackslash}m{2.2cm}|>{\centering\arraybackslash}m{2.2cm}|>{\centering\arraybackslash}m{2.2cm}|>{\centering\arraybackslash}m{2.2cm}|}
\hline
\textbf{Exercice 1} & \textbf{Exercice 2} & \textbf{Exercice 3} & \textbf{Exercice 4} & \textbf{Exercice 5} & \textbf{Total} \\
\hline
4 pts & 4 pts & 4 pts & 4 pts & 4 pts & 20 pts \\
\hline
\end{tabular}
\end{center}

\section*{Exercice 1 -- Mouvement et vitesse \hfill /4}

Un cycliste parcourt une distance de \SI{3,6}{km} en \SI{12}{min}.

\textbf{Rappels utiles}
\begin{itemize}[leftmargin=1.2cm]
    \item La vitesse moyenne se calcule par :
    \[
    v=\frac{d}{t}
    \]
    \item Pour utiliser cette formule correctement, il faut exprimer les grandeurs dans des unités cohérentes.
    \item \SI{1}{h} = \SI{60}{min}.
\end{itemize}

\begin{enumerate}
    \item Convertir la durée du trajet en heure.
    \item Calculer la vitesse moyenne du cycliste en \si{km.h^{-1}}.
    \item Convertir cette vitesse en \si{m.s^{-1}}.
    \item Le mouvement du cycliste est-il nécessairement uniforme ? Expliquer.
\end{enumerate}

\textbf{Barème indicatif :} 1 pt ; 1,5 pt ; 1 pt ; 0,5 pt.

\section*{Exercice 2 -- Électricité \hfill /4}

On étudie un appareil électrique portant les indications suivantes :

\begin{center}
\framebox{
\begin{minipage}{0.65\textwidth}
\centering
\textbf{Lampe LED}\\
Tension nominale : \SI{6}{V}\\
Intensité nominale : \SI{0,30}{A}
\end{minipage}}
\end{center}

\textbf{Rappels utiles}
\begin{itemize}[leftmargin=1.2cm]
    \item La puissance électrique se calcule par :
    \[
    P=U\times I
    \]
    \item Avec \(P\) en watt (W), \(U\) en volt (V) et \(I\) en ampère (A).
\end{itemize}

\begin{enumerate}
    \item Identifier la tension nominale et l’intensité nominale de la lampe.
    \item Calculer la puissance électrique de cette lampe.
    \item On alimente cette lampe pendant \SI{10}{min}. Calculer l’énergie consommée en watt-heure (\si{Wh}).
    \item Citer un dipôle de protection que l’on trouve dans un circuit domestique et préciser son rôle.
\end{enumerate}

\textbf{Barème indicatif :} 1 pt ; 1 pt ; 1,5 pt ; 0,5 pt.

\section*{Exercice 3 -- Signaux et propagation du son \hfill /4}

Lors d’un orage, une élève voit un éclair puis entend le tonnerre \SI{5}{s} plus tard.

On donne la vitesse du son dans l’air :
\[
v_{\text{son}} \approx \SI{340}{m.s^{-1}}
\]

\textbf{Rappels utiles}
\begin{itemize}[leftmargin=1.2cm]
    \item Le son met un certain temps à se propager.
    \item La lumière se propage beaucoup plus vite que le son.
    \item On peut utiliser la relation :
    \[
    d=v\times t
    \]
\end{itemize}

\begin{enumerate}
    \item Expliquer pourquoi on voit l’éclair avant d’entendre le tonnerre.
    \item Calculer la distance approximative entre l’orage et l’élève.
    \item Exprimer le résultat en kilomètre.
    \item Dire si le son a besoin d’un milieu matériel pour se propager. Justifier brièvement.
\end{enumerate}

\textbf{Barème indicatif :} 1 pt ; 1,5 pt ; 0,5 pt ; 1 pt.

\section*{Exercice 4 -- Transformation chimique \hfill /4}

On fait réagir du carbone avec du dioxygène. Il se forme du dioxyde de carbone.

\textbf{Rappels utiles}
\begin{itemize}[leftmargin=1.2cm]
    \item Lors d’une transformation chimique, les réactifs disparaissent et les produits apparaissent.
    \item Les atomes se conservent.
    \item Le dioxygène est noté \(O_2\), le dioxyde de carbone est noté \(CO_2\).
\end{itemize}

\begin{enumerate}
    \item Nommer les réactifs et le produit de cette transformation chimique.
    \item Écrire l’équation-bilan ajustée :
    \[
    \cdots + \cdots \rightarrow \cdots
    \]
    \item Expliquer pourquoi on dit qu’il y a conservation des atomes.
    \item Citer un test permettant d’identifier le dioxyde de carbone.
\end{enumerate}

\textbf{Barème indicatif :} 1 pt ; 1 pt ; 1 pt ; 1 pt.

\section*{Exercice 5 -- Masse volumique et identification d’un matériau \hfill /4}

On mesure la masse d’un objet métallique : 
\[
m=\SI{54}{g}
\]
et son volume :
\[
V=\SI{20}{cm^3}.
\]

On rappelle que la masse volumique se calcule par :
\[
\rho=\frac{m}{V}
\]

On donne le tableau suivant :

\begin{center}
\begin{tabular}{|c|c|}
\hline
\textbf{Matériau} & \textbf{Masse volumique (\si{g.cm^{-3}})} \\
\hline
Aluminium & 2,7 \\
Fer & 7,8 \\
Cuivre & 8,9 \\
Plastique & 1,2 \\
\hline
\end{tabular}
\end{center}

\begin{enumerate}
    \item Calculer la masse volumique de l’objet.
    \item En utilisant le tableau, proposer la nature la plus probable du matériau.
    \item L’objet flotte-t-il dans l’eau ? On rappelle que la masse volumique de l’eau vaut environ \SI{1,0}{g.cm^{-3}}.
    \item Expliquer pourquoi la masse volumique peut aider à identifier une substance.
\end{enumerate}

\textbf{Barème indicatif :} 1,5 pt ; 1 pt ; 0,5 pt ; 1 pt.

\vspace{0.5cm}

\begin{center}
\textbf{Fin du sujet -- Toute réponse doit être justifiée avec soin.}
\end{center}

\end{document}